\documentclass[11pt,a4paper]{article}

\usepackage[utf8]{inputenc}
\usepackage[T1]{fontenc}
\usepackage{lmodern}
\usepackage{amsmath,amssymb,amsfonts}
\usepackage{booktabs}
\usepackage{geometry}
\usepackage{hyperref}
\usepackage{graphicx}
\usepackage{physics}

\geometry{margin=2.5cm}

\title{
Paper 12 --- SPARC Rotation Curves from Local Entanglement Profiles\\
\large Bottom-Up Quantum Gravity
}

\author{Farid Hamdad}
\date{2026}

\begin{document}

\maketitle

\begin{abstract}
We test the effective propagator law derived in Paper 11 of the Bottom-Up
Quantum Gravity (BuP) program on SPARC galactic rotation curves. In BuP,
gravity is not introduced as a fundamental force or as a dark matter halo, but
emerges from the entanglement structure of a finite quantum state. Paper 11
showed that the effective power-law exponent of the entanglement propagator is
controlled not only by the spectral dimension \(d_s\), but by the pair
\((d_s,d_w)\), where \(d_w\) is the walk dimension:
\[
\alpha_{\rm eff}
=
\frac{2d_s}{d_w}+d_w-4.
\]
Paper 12 applies this relation to baryonic SPARC disks by constructing local
entanglement graphs, measuring radial profiles \(d_s(r)\), \(d_w(r)\), deriving
\(\alpha_{\rm eff}(r)\), and projecting the result into a correction to the
circular velocity.

The v3.4 corrected model improves 17 out of 18 galaxies in an extended
best-window sample, with 8 galaxies reaching \(\chi^2_{\rm red}<5\). Including
the targeted NGC5055 compacity correction yields an effective 19-galaxy
classification with 18 improved galaxies, 8 galaxies below
\(\chi^2_{\rm red}<5\), and 12 below \(\chi^2_{\rm red}<10\). The results do
not yet constitute a final universal SPARC model. Instead, they show that
galaxies split into classes of entanglement profiles: regular transitions,
extended disks, internally complex systems, and diffuse early sub-Newtonian
regimes.
\end{abstract}

\section{Introduction}

The SPARC rotation-curve problem is usually formulated as a discrepancy between
the observed circular velocity \(V_{\rm obs}(r)\) and the baryonic prediction
\(V_{\rm bar}(r)\). In standard dark-matter phenomenology, the missing
contribution is modeled by a halo. In the Bottom-Up Quantum Gravity approach,
the additional effective contribution is instead interpreted as a geometric
response of the entanglement structure.

The central idea of BuP is that geometry is not fundamental. An entanglement
matrix \(W_{ij}\) defines a graph, the graph defines a Laplacian
\(L_{\rm ent}\), and the pseudo-inverse or propagator of this Laplacian defines
an effective gravitational response. The relevant chain is
\[
W_{ij}
\rightarrow
L_{\rm ent}
\rightarrow
(d_s,d_w)
\rightarrow
\alpha_{\rm eff}
\rightarrow
\Phi(r)
\rightarrow
V(r).
\]

Earlier SPARC studies in the BuP program used a scalar effective dimension
\(d(r)\). These runs already indicated that SPARC galaxies are not described by
a single universal regime. The old profiles \(d_{\rm emp}(r)\) and
\(d_{\rm fit}(r)\) separated galaxies into groups such as
\texttt{spiral\_massive}, \texttt{spiral\_intermediate}, and
\texttt{dwarf\_gas}. Paper 12 reformulates this empirical separation in terms
of local entanglement profiles:
\[
d(r)
\quad\longrightarrow\quad
(d_s(r),d_w(r))
\quad\longrightarrow\quad
\alpha_{\rm eff}(r).
\]

\section{Link with Paper 11}

Paper 11 established the key theoretical input of the present work. For a graph
whose heat kernel is characterized by a spectral dimension \(d_s\) and a walk
dimension \(d_w\), the effective propagator exponent is not simply
\(d_s-2\). The BuP effective exponent is
\[
\boxed{
\alpha_{\rm eff}
=
\frac{2d_s}{d_w}+d_w-4.
}
\]

When \(d_w=2\), the Brownian limit is recovered:
\[
\alpha_{\rm eff}=d_s-2.
\]
In the cosmological quasi-three-dimensional regime, this relation approaches a
nearly Newtonian exponent. However, Paper 12 emphasizes that galaxies need not
live in the same entanglement regime as the large-scale cosmological graph. A
galactic disk can generate its own local pair \((d_s(r),d_w(r))\), and hence a
radially dependent exponent \(\alpha_{\rm eff}(r)\).

This is the conceptual bridge between Paper 11 and Paper 12:
\[
\boxed{
\text{Paper 11 derives the propagator law; Paper 12 tests it on SPARC.}
}
\]

\section{Baryonic entanglement graph}

For each galaxy, the observed baryonic disk is converted into a two-dimensional
set of graph nodes. The graph is weighted by a baryonic correlation scale
\(\lambda_{\rm corr}\), a local connectivity \(k\), and the baryonic mass
distribution inferred from the SPARC rotation decomposition.

The default graph parameters used in the present tests are
\[
N_{\rm nodes}=600,
\qquad
k=12,
\qquad
\lambda_{\rm corr}=3.0\,{\rm kpc}.
\]

From the resulting graph, local windows are selected in radius. For each radial
window, we compute:
\[
d_s(r)
\]
from heat-kernel scaling, and
\[
d_w(r)
\]
from mean-square displacement scaling. These define
\[
\alpha_{\rm eff}(r)
=
\frac{2d_s(r)}{d_w(r)}
+
d_w(r)
-
4.
\]

\section{Projection to rotation curves}

The present version does not yet compute the full non-local propagator for the
galaxy. Instead, it uses the profile \(\alpha_{\rm eff}(r)\) to determine the
onset of an external BuP contribution to \(V^2\). The velocity model is
\[
V^2_{\rm model}(r)
=
V^2_{\rm bar}(r)
+
A\,S(r),
\]
where
\[
S(r)
=
\frac{1}{1+\exp[-(r-r_t)/\Delta r]}.
\]

The transition radius \(r_t\) is not chosen freely in the v3.3/v3.4 models. It
is derived from the local profile \(\alpha_{\rm eff}(r)\). The relative trigger
is
\[
\alpha_{\rm eff}(r_t)
=
q\,\alpha_{\max},
\qquad
q=0.97.
\]

The transition width is fixed to
\[
\Delta r=2.5\,{\rm kpc}.
\]

\section{v3.4 corrections}

The v3.4 version adds two corrections motivated by earlier SPARC runs.

\subsection{Adaptive local extraction window}

A fixed radial window is not optimal for all galaxies. Extended regular disks
can require a larger local extraction scale to stabilize
\((d_s,d_w,\alpha_{\rm eff})\). We therefore scan
\[
w\in\{5,8,10\}\,{\rm kpc}
\]
and select the best local extraction window for the class analysis. This is not
a dynamical halo parameter; it is an extraction scale for the local entanglement
profile.

\subsection{Compacity correction}

For galaxies with internally complex baryonic structure, the profile
\(\alpha_{\rm eff}(r)\) can develop a high central peak. The old hierarchical
SPARC runs already indicated that such galaxies require a compacity correction.
We define
\[
\eta
=
\frac{v_{b,\max}-v_{b,\rm med}}{v_{b,\rm med}}.
\]

For high-peak profiles,
\[
\alpha_{\max}>8,
\]
we correct the transition radius by
\[
\boxed{
r_t^{\rm corr}
=
r_t^{(0)}(1-0.75\eta).
}
\]

This correction advances the BuP transition in compact or internally complex
systems.

\section{Canonical cases}

\subsection{NGC3198: regular transition}

NGC3198 is the canonical Paper 12 case. With
\[
w=5\,{\rm kpc},
\qquad
q=0.97,
\]
the profile gives
\[
\alpha_{\max}=4.965,
\qquad
\alpha_{\rm med}=0.369,
\]
and
\[
r_t=5.18\,{\rm kpc}.
\]
The resulting fit is
\[
\chi^2_{\rm red}=0.994,
\]
compared with
\[
\chi^2_{\rm red,bar}=759.05.
\]

\subsection{NGC2841: extended standard disk}

Earlier BuP runs described NGC2841 as a standard stable galaxy. A fixed
\(w=5\,{\rm kpc}\) window in Paper 12 v3.3 degraded the fit:
\[
\chi^2_{\rm red}=42.0.
\]
Using a larger extraction window,
\[
w=10\,{\rm kpc},
\]
regularizes the local profile:
\[
\alpha_{\max}=1.87,
\qquad
\alpha_{\rm med}=0.62,
\]
and yields
\[
\chi^2_{\rm red}=4.59.
\]
This shows that NGC2841 is not intrinsically complex; it requires a local
window adapted to its radial extension.

\subsection{NGC5055: compacity-corrected complex galaxy}

NGC5055 is the clearest complex case. With \(w=5\,{\rm kpc}\), its local
profile contains a high peak:
\[
\alpha_{\max}=12.60.
\]
Without compacity correction, the transition is too late:
\[
r_t^{(0)}=16.35\,{\rm kpc},
\qquad
\chi^2_{\rm red}=111.3.
\]
The baryonic compacity is
\[
\eta=0.337.
\]
Applying
\[
r_t^{\rm corr}=r_t^{(0)}(1-0.75\eta)
\]
gives
\[
r_t^{\rm corr}=12.22\,{\rm kpc},
\]
and reduces the fit to
\[
\chi^2_{\rm red}=5.53.
\]

Wider windows erase the high-\(\alpha\) peak and prevent the compacity
correction from activating. Therefore the peak is interpreted as a physical
signature of internal complexity, not as numerical noise.

\subsection{NGC2403: diffuse early sub-Newtonian regime}

NGC2403 remains the main resistant case. Earlier BuP versions already gave
elevated values \(\chi^2_{\rm red}\sim 11-15\), indicating that this galaxy is
not a new failure mode of Paper 12. In the present language, its profile is
strongly sub-Newtonian almost from the center:
\[
\alpha_{\rm med}=0.0818,
\qquad
r_t=0.52\,{\rm kpc}.
\]
A single external sigmoid is therefore not appropriate. NGC2403 defines a
third regime: diffuse early sub-Newtonian transition, probably requiring a
non-sigmoidal or multi-scale projection.

\section{Extended sample}

The v3.4 best-window scan was run on an extended sample of 18 galaxies. The
summary is:

\begin{table}[h]
\centering
\begin{tabular}{lc}
\toprule
Metric & Value \\
\midrule
Number of galaxies & 18 \\
Improved vs baryon-only & 17 \\
\(\chi^2_{\rm red}<2\) & 4 \\
\(\chi^2_{\rm red}<5\) & 8 \\
\(\chi^2_{\rm red}<10\) & 11 \\
Median \(\chi^2_{\rm red}\) & 6.80 \\
Mean \(\chi^2_{\rm red}\) & 14.10 \\
\bottomrule
\end{tabular}
\end{table}

The best cases are:

\begin{table}[h]
\centering
\begin{tabular}{lccc}
\toprule
Galaxy & \(w\) (kpc) & \(\alpha_{\max}\) & \(\chi^2_{\rm red}\) \\
\midrule
NGC3198 & 5 & 4.97 & 0.994 \\
NGC3769 & 5 & 5.37 & 1.08 \\
NGC5005 & 10 & 0.82 & 1.43 \\
NGC3521 & 10 & 1.70 & 1.82 \\
NGC7793 & 8 & 2.15 & 2.55 \\
NGC2998 & 8 & 4.85 & 2.62 \\
NGC0024 & 10 & 1.42 & 3.95 \\
NGC2841 & 10 & 1.87 & 4.59 \\
\bottomrule
\end{tabular}
\end{table}

Including the targeted NGC5055 compacity-corrected result gives an effective
19-galaxy classification:

\begin{table}[h]
\centering
\begin{tabular}{lc}
\toprule
Metric & Value \\
\midrule
Number of galaxies & 19 \\
Improved vs baryon-only & 18 \\
\(\chi^2_{\rm red}<2\) & 4 \\
\(\chi^2_{\rm red}<5\) & 8 \\
\(\chi^2_{\rm red}<10\) & 12 \\
Median \(\chi^2_{\rm red}\) & 6.61 \\
Mean \(\chi^2_{\rm red}\) & 13.65 \\
\bottomrule
\end{tabular}
\end{table}

\section{Classification}

The results suggest four BuP regimes.

\subsection{Class I-A: regular transition}

These galaxies are well described by a short local extraction window and a
regular relative trigger. NGC3198 is the canonical example.

\subsection{Class I-B: extended standard disks}

These galaxies are standard BuP systems but require a larger local extraction
window. NGC2841 is the canonical example.

\subsection{Class II: internally complex systems}

These galaxies require a structural or compacity correction. NGC5055 is the
canonical example. The high \(\alpha_{\max}\) peak signals internal complexity,
and the transition radius must be advanced by the \(\eta\) correction.

\subsection{Class III: diffuse early sub-Newtonian systems}

These galaxies cannot be described by a single external sigmoid. NGC2403 is
the canonical example. Future versions must replace the sigmoid by a diffuse
or multi-scale projection.

\section{Discussion}

Paper 12 does not claim that the present v3.4 model is a final universal
description of SPARC. Its main contribution is more fundamental: it shows that
the effective propagator law from Paper 11 has observable consequences in
galactic rotation curves, and that SPARC galaxies organize into classes of
local entanglement profiles.

This also explains why earlier BuP runs found good fits for some galaxies and
failures for others. The failures are not arbitrary. They correspond to
different regimes of \((d_s(r),d_w(r),\alpha_{\rm eff}(r))\), requiring
different projections.

\section{Conclusion}

Paper 11 derived the BuP effective propagator exponent
\[
\alpha_{\rm eff}
=
\frac{2d_s}{d_w}+d_w-4.
\]
Paper 12 applied this exponent to SPARC rotation curves by measuring local
entanglement profiles in baryonic graphs. The v3.4 results show that the method
improves almost all tested galaxies and gives good fits for a significant
subsample.

The main conclusion is:
\[
\boxed{
\text{SPARC galaxies are organized by classes of entanglement profiles.}
}
\]

The next step is to replace the single sigmoid projection by class-dependent
non-local propagator reconstructions, especially for diffuse systems such as
NGC2403.

\end{document}
